\documentclass[12pt]{article}
\usepackage{amsmath}
\usepackage{amsfonts}
\usepackage{float}
\usepackage{fancyhdr}
\usepackage{graphicx}
\usepackage{setspace}
\usepackage[colorlinks=true,linkcolor=blue, citecolor=red]{hyperref}
\usepackage{url}
\usepackage[top=.75in, left=.75in, right=.75in, bottom=1in]{geometry}

% For algorithm
\usepackage{algorithm}
\usepackage{algpseudocode}

% ============ CODE ============
\usepackage{listings}
\usepackage{xcolor}
\definecolor{codegreen}{rgb}{0,0.6,0}
\definecolor{codegray}{rgb}{0.5,0.5,0.5}
\definecolor{codepurple}{rgb}{0.58,0,0.82}
\definecolor{backcolour}{rgb}{0.95,0.95,0.92}

% Styling for the code.
\lstdefinestyle{mystyle}{
    backgroundcolor=\color{backcolour},   
    commentstyle=\color{codegreen},
    keywordstyle=\color{magenta},
    numberstyle=\tiny\color{codegray},
    stringstyle=\color{codepurple},
    basicstyle=\ttfamily\footnotesize,
    breakatwhitespace=false,         
    breaklines=true,                 
    captionpos=b,                    
    keepspaces=true,                 
    numbers=left,                    
    numbersep=5pt,                  
    showspaces=false,                
    showstringspaces=false,
    showtabs=false,                  
    tabsize=2
}
\lstset{style=mystyle}

% Table packages and Unicode / Font support
\usepackage{array}
\usepackage{longtable}
\usepackage{enumitem}
\usepackage{booktabs}
\usepackage[utf8]{inputenc}
\DeclareUnicodeCharacter{2264}{\ensuremath{\le}}
\DeclareUnicodeCharacter{2194}{\ensuremath{\leftrightarrow}}
\usepackage[T1]{fontenc}
\usepackage{lmodern}
\graphicspath{{../../../../templates/}{../../prototype/}{figures/}{./}}
\setlength{\emergencystretch}{3em}
\setlist[enumerate]{label=\arabic*.,leftmargin=*}
\setlist[itemize]{leftmargin=*}

% ============ HEADER AND FOOTER ============
\setlength{\headheight}{29.43912pt}

\pagestyle{fancy}
\fancyfoot{}
\fancyfoot[R]{Page \thepage}

\setlength{\parindent}{0pt}
\renewcommand{\baselinestretch}{1.5}

\newcommand{\coursename}{Introduction of Software Engineering}
\newcommand{\reportname}{Studitopia}
\newcommand{\reporttitle}{\href{https://github.com/RoundFloor/Studitopia.git}{Studitopia}}

\newcommand{\studentname}{%
Nguyen Thanh An (24125024)\\
Nguyen Chi Bao (24125025)\\
Tran Nguyen Khanh Binh (24125026)\\
Tran Canh Anh Tuan (24125048)\\
Nguyen Le Quoc Huy (24125100)
}
\newcommand{\teachername}{Ho Tuan Thanh, Nguyen Thi Minh Tuyen}

% Header
\lhead{\reporttitle}
\rhead{
University of Science - VNUHCM\\
\coursename
}

% Footer
\newcommand{\leftfooter}{}
\lfoot{\leftfooter}

% ============ DOCUMENT ============
\begin{document}

\pagenumbering{roman}
\begin{titlepage}
\newcommand{\HRule}{\rule{\linewidth}{0.5mm}}
\centering

\textsc{\LARGE VNU-HCM}\\[0.5cm]
\textsc{\Large University of Science}\\[0.5cm]
\textsc{\large Faculty of Information Technology}\\[0.5cm]
\textsc{Computer Science}\\[0.5cm]

\HRule \\[0.4cm]
{ 
\huge{\bfseries{\reporttitle}}\\[0.5cm]
\large{\bfseries{UC-14 --- Manage Privacy}}
}\\[0.4cm]
\HRule \\[0.5cm]

\textbf{\large Subject: \coursename}\\[0.5cm]

\begin{minipage}[t]{0.55\textwidth}
\begin{flushleft} \large
\emph{Students:}\\
\studentname
\end{flushleft}
\end{minipage}
~
\begin{minipage}[t]{0.4\textwidth}
\begin{flushright} \large
\emph{Supervising Teachers:} \\
\teachername
\end{flushright}
\end{minipage}\\[1cm]

{\large \today}\\[1cm]

\includegraphics[scale=.20]{img/hcmus-logo.png}\\[1cm] 

\vfill
\end{titlepage}	

\tableofcontents
\pagebreak

\pagenumbering{arabic}
\setcounter{page}{1}

\hypertarget{sec:uc-14-manage-privacy}{}
\addcontentsline{toc}{section}{UC-14 --- Manage Privacy}
\section*{UC-14 --- Manage Privacy}
\par\medskip

\textbf{Version:} \texttt{1.0}
\par\medskip

\textbf{Performed by:} [Tran Canh Anh Tuan] | \textbf{Edited by:} [Tran Canh Anh Tuan]
\par\medskip

\begingroup
\small
\renewcommand{\baselinestretch}{1.0}\selectfont
\renewcommand{\arraystretch}{1.15}
\setlength{\tabcolsep}{6pt}
\begin{longtable}{|>{\raggedright\arraybackslash}p{0.10\textwidth}|>{\raggedright\arraybackslash}p{0.83\textwidth}|}
\hline
\textbf{Field} & \textbf{Value} \\ \hline
\endhead
Use Case ID & UC-14 \\ \hline
Name & Manage Privacy \\ \hline
Primary Actor & User \\ \hline
Supporting Actors & None \\ \hline
Feature Group & Profile Management \\ \hline
Priority & Not explicitly stated \\ \hline
Related User Stories & US-10 \\ \hline
Related Requirements & FR-10.1, FR-10.2, FR-10.3, FR-10.4, FR-39.1, FR-39.2, FR-39.3, FR-39.4, FR-39.5; NFR-PRIV-1, NFR-PRIV-2, NFR-PRIV-3, NFR-PRIV-4, NFR-PRIV-5, NFR-PRIV-6, NFR-PRIV-7, NFR-SEC-1, NFR-SEC-3, NFR-SEC-4, NFR-SEC-6, NFR-USE-4 \\ \hline
Source Documents & \texttt{proposals/userstory.md} (§3.3); \texttt{proposals/functional\_requirement.md} (US-10, US-39); \texttt{proposals/settings.md}; \texttt{proposals/friend\_profile.md}; \texttt{proposals/overview.md}; \texttt{proposals/feature\_flow.md}; \texttt{proposals/screen\_flow.md}; \texttt{proposals/API.md}; \texttt{proposals/non\_functional\_requirement.md} \\ \hline
\end{longtable}
\endgroup
\par\medskip

\hypertarget{sec:part-a-use-case-model}{}
\addcontentsline{toc}{subsection}{Part A --- Use Case Model}
\subsection*{Part A --- Use Case Model}
\par\medskip

\hypertarget{sec:a1-focused-use-case-diagram}{}
\addcontentsline{toc}{subsubsection}{A1. Focused Use Case Diagram}
\subsubsection*{A1. Focused Use Case Diagram}
\par\medskip

\begin{center}
\includegraphics[width=0.85\textwidth,height=0.75\textheight,keepaspectratio]{"figures/UC-14-manage-privacy_1.png"}
\end{center}
\par\medskip

\textbf{Diagram notation.} Person shape = human actor; cloud = supporting system; blue = selected use case; solid line = association; labelled solid arrow = generalization; dashed arrow = include or extend.
\par\medskip

\hypertarget{sec:a2-actor-register}{}
\addcontentsline{toc}{subsubsection}{A2. Actor Register}
\subsubsection*{A2. Actor Register}
\par\medskip

\begingroup
\small
\renewcommand{\baselinestretch}{1.0}\selectfont
\renewcommand{\arraystretch}{1.15}
\setlength{\tabcolsep}{6pt}
\begin{longtable}{|>{\raggedright\arraybackslash}p{0.09\textwidth}|>{\raggedright\arraybackslash}p{0.09\textwidth}|>{\raggedright\arraybackslash}p{0.74\textwidth}|}
\hline
\textbf{Actor} & \textbf{Type} & \textbf{Description} \\ \hline
\endhead
User & Primary & A platform user who controls the visibility of personal information and limits communication to protect privacy. \\ \hline
\end{longtable}
\endgroup
\par\medskip

\hypertarget{sec:a3-relationships}{}
\addcontentsline{toc}{subsubsection}{A3. Relationships}
\subsubsection*{A3. Relationships}
\par\medskip

\hypertarget{sec:a3-1-associations}{}
\paragraph*{A3.1 Associations}
\par\medskip

\begingroup
\small
\renewcommand{\baselinestretch}{1.0}\selectfont
\renewcommand{\arraystretch}{1.15}
\setlength{\tabcolsep}{6pt}
\begin{longtable}{|>{\raggedright\arraybackslash}p{0.09\textwidth}|>{\raggedright\arraybackslash}p{0.11\textwidth}|>{\raggedright\arraybackslash}p{0.09\textwidth}|>{\raggedright\arraybackslash}p{0.62\textwidth}|}
\hline
\textbf{Actor} & \textbf{Use Case} & \textbf{Type} & \textbf{Description} \\ \hline
\endhead
User & UC-14 Manage Privacy & Primary & Opens \textbf{Privacy} or \textbf{Communication} in \textbf{Settings} and saves visibility and interaction restrictions. \\ \hline
\end{longtable}
\endgroup
\par\medskip

\hypertarget{sec:a3-2-generalization-and-specialization}{}
\paragraph*{A3.2 Generalization and Specialization}
\par\medskip

\begingroup
\small
\renewcommand{\baselinestretch}{1.0}\selectfont
\renewcommand{\arraystretch}{1.15}
\setlength{\tabcolsep}{6pt}
\begin{longtable}{|>{\raggedright\arraybackslash}p{0.09\textwidth}|>{\raggedright\arraybackslash}p{0.09\textwidth}|>{\raggedright\arraybackslash}p{0.12\textwidth}|>{\raggedright\arraybackslash}p{0.61\textwidth}|}
\hline
\textbf{Child} & \textbf{Parent} & \textbf{Type} & \textbf{Inherited Behavior or Association} \\ \hline
\endhead
None & None & No relationship & No actor generalization or specialization is named for US-10 in the proposals. \\ \hline
\end{longtable}
\endgroup
\par\medskip

\hypertarget{sec:a3-3-include}{}
\paragraph*{A3.3 Include}
\par\medskip

\begingroup
\small
\renewcommand{\baselinestretch}{1.0}\selectfont
\renewcommand{\arraystretch}{1.15}
\setlength{\tabcolsep}{6pt}
\begin{longtable}{|>{\raggedright\arraybackslash}p{0.14\textwidth}|>{\raggedright\arraybackslash}p{0.18\textwidth}|>{\raggedright\arraybackslash}p{0.60\textwidth}|}
\hline
\textbf{Base Use Case} & \textbf{Included Use Case} & \textbf{Meaning} \\ \hline
\endhead
None & None & No included use case is named for US-10 in the proposals. \\ \hline
\end{longtable}
\endgroup
\par\medskip

\hypertarget{sec:a3-4-extend}{}
\paragraph*{A3.4 Extend}
\par\medskip

\begingroup
\small
\renewcommand{\baselinestretch}{1.0}\selectfont
\renewcommand{\arraystretch}{1.15}
\setlength{\tabcolsep}{6pt}
\begin{longtable}{|>{\raggedright\arraybackslash}p{0.14\textwidth}|>{\raggedright\arraybackslash}p{0.10\textwidth}|>{\raggedright\arraybackslash}p{0.21\textwidth}|>{\raggedright\arraybackslash}p{0.46\textwidth}|}
\hline
\textbf{Extending Use Case} & \textbf{Base Use Case} & \textbf{Condition} & \textbf{Added Behavior} \\ \hline
\endhead
None & None & Not specified in proposals & No extending use case is named for US-10 in the proposals. \\ \hline
\end{longtable}
\endgroup
\par\medskip

\hypertarget{sec:a4-model-notes}{}
\addcontentsline{toc}{subsubsection}{A4. Model Notes}
\subsubsection*{A4. Model Notes}
\par\medskip

\begin{enumerate}
\item Scope includes visibility, discoverability, online status, and communication boundaries grouped under \textbf{Privacy} and \textbf{Communication} in \texttt{settings.md}.
\item Overall profile visibility values \textbf{Public / Private} come from \texttt{settings.md} and NFR-PRIV-1; field-level values \textbf{Visible / Hidden} come from AC-10.1, FR-10.1, and the profile settings flow.
\item No dedicated privacy API endpoint is documented in \texttt{API.md}; all route names in \texttt{settings.md} are recommendations.
\end{enumerate}
\par\medskip

\hypertarget{sec:part-b-use-case-specification}{}
\addcontentsline{toc}{subsection}{Part B --- Use Case Specification}
\subsection*{Part B --- Use Case Specification}
\par\medskip

\hypertarget{sec:b1-brief-description}{}
\addcontentsline{toc}{subsubsection}{B1. Brief Description}
\subsubsection*{B1. Brief Description}
\par\medskip

A user controls the visibility of personal information so that they can protect their privacy while using the platform. Controls cover overall profile visibility, field-level visibility for academic, study-history, habit, and contact information, searchability, online status, and last-active time.
\par\medskip

The user also controls who may communicate with them through direct-message restrictions, friend-request and room-invitation controls, and muted or blocked member lists. Privacy settings must be respected wherever user information appears in posts, comments, and member lists.
\par\medskip

\hypertarget{sec:b2-preconditions}{}
\addcontentsline{toc}{subsubsection}{B2. Preconditions}
\subsubsection*{B2. Preconditions}
\par\medskip

\begin{enumerate}
\item The user is authenticated before accessing protected privacy settings (NFR-SEC-1).
\item The user has a profile whose fields can be made \textbf{Visible} or \textbf{Hidden} (FR-10.1, AC-10.1).
\item The user has opened \textbf{Settings} and selected \textbf{Privacy} or \textbf{Communication}.
\item Supported member identifiers are available for managing muted and blocked lists (FR-39.3).
\end{enumerate}
\par\medskip

\hypertarget{sec:b3-postconditions}{}
\addcontentsline{toc}{subsubsection}{B3. Postconditions}
\subsubsection*{B3. Postconditions}
\par\medskip

\begin{enumerate}
\item \textbf{Success:} Saved privacy and communication values are stored and applied consistently to posts, comments, member lists, and search results (FR-10.4, AC-10.4, NFR-PRIV-7).
\item \textbf{Success:} Disallowed messages, friend requests, or room invitations are prevented by backend enforcement (FR-39.5, NFR-SEC-6).
\item \textbf{Failure:} Unsupported values are not saved and the previously saved privacy settings remain in effect.
\end{enumerate}
\par\medskip

\hypertarget{sec:b4-main-success-scenario-basic-flow}{}
\addcontentsline{toc}{subsubsection}{B4. Main Success Scenario (Basic Flow)}
\subsubsection*{B4. Main Success Scenario (Basic Flow)}
\par\medskip

\begin{enumerate}
\item The user opens \textbf{Settings} and selects \textbf{Privacy}.
\item The system displays \texttt{PrivacySettingsForm} with \texttt{ProfileVisibilitySelect} (\textbf{Public / Private}), field-level \texttt{AcademicInfoVisibilitySelect}, \texttt{ContactInfoVisibilitySelect}, \texttt{StudyHistoryVisibilitySelect}, and \texttt{HabitsVisibilitySelect} (\textbf{Visible / Hidden}), plus \texttt{OnlineStatusToggle}, \texttt{LastActiveVisibilityToggle}, \texttt{PublicDiscoveryOptOutToggle}, and \texttt{StudyPartnerSearchOptOutToggle}.
\item The user chooses the overall profile visibility, field-level values, and search or status toggles.
\item The user opens \textbf{Communication} and sets \texttt{DirectMessageRestrictionSelect} (\textbf{Everyone / Confirmed friends only}), \texttt{FriendRequestPermissionSelect} (\textbf{Allow / Reject incoming friend requests}), and \texttt{RoomInvitationPermissionSelect} (\textbf{Allow / Reject room invitations}).
\item The user manages \texttt{MutedUsersList} and \texttt{BlockedUsersList} with \textbf{Mute}, \textbf{Unmute}, \textbf{Block}, and \textbf{Unblock}.
\item The system shows the note \textbf{Applied to posts, comments, member lists, and search results}.
\item The user selects \textbf{Save Privacy} or \textbf{Save Communication}.
\item Backend stores the values and applies them to relevant reads and interactions.
\item The system displays a success toast.
\end{enumerate}
\par\medskip

\hypertarget{sec:b5-alternative-flows}{}
\addcontentsline{toc}{subsubsection}{B5. Alternative Flows}
\subsubsection*{B5. Alternative Flows}
\par\medskip

No distinct alternative flow is stated in US-10; the following source-backed variants are recorded.
\par\medskip

\hypertarget{sec:a1-communication-only-change}{}
\paragraph*{A1. Communication-only change}
\par\medskip

At basic flow step 1, the user opens \textbf{Communication} directly and changes only message, request, invitation, mute, or block controls before step 7.
\par\medskip

\hypertarget{sec:a2-public-discovery-or-study-partner-search-opt-out}{}
\paragraph*{A2. Public discovery or study-partner search opt-out}
\par\medskip

At basic flow step 3, the user enables \texttt{PublicDiscoveryOptOutToggle} or \texttt{StudyPartnerSearchOptOutToggle} to opt out of the corresponding search results (FR-10.2, AC-10.2).
\par\medskip

\hypertarget{sec:a3-online-status-and-last-active-change}{}
\paragraph*{A3. Online status and last-active change}
\par\medskip

At basic flow step 3, the user hides or shows online status and last-active time (FR-10.3, AC-10.3, NFR-PRIV-4).
\par\medskip

\hypertarget{sec:b6-exception-flows}{}
\addcontentsline{toc}{subsubsection}{B6. Exception Flows}
\subsubsection*{B6. Exception Flows}
\par\medskip

No exception flow is specified in US-10 or FR-10.1 to FR-10.4. No status code or error code for unsupported privacy values is specified in proposals.
\par\medskip

\hypertarget{sec:b7-business-rules}{}
\addcontentsline{toc}{subsubsection}{B7. Business Rules}
\subsubsection*{B7. Business Rules}
\par\medskip

\begin{enumerate}
\item Overall profile visibility uses \textbf{Public / Private} and provides at least two visibility levels (NFR-PRIV-1).
\item Field-level visibility uses \textbf{Visible / Hidden} for study history, habits, and contact information (FR-10.1, AC-10.1).
\item University, major, and academic year visibility is controllable (NFR-PRIV-3).
\item Online status and last-active time can be shown or hidden (FR-10.3, AC-10.3, NFR-PRIV-4).
\item Users can opt out of public discovery and study-partner search results (FR-10.2, AC-10.2).
\item Privacy settings are enforced in posts, comments, member lists, and search results (FR-10.4, AC-10.4, NFR-PRIV-7).
\item Direct messages are restricted to \textbf{Everyone / Confirmed friends only} (FR-39.4, AC-39.2, NFR-PRIV-6).
\item Friend requests offer \textbf{Allow / Reject incoming friend requests} and room invitations offer \textbf{Allow / Reject room invitations} (NFR-PRIV-2, NFR-PRIV-6).
\item Muted and blocked members are managed with \textbf{Mute}, \textbf{Unmute}, \textbf{Block}, and \textbf{Unblock} (FR-39.3, AC-39.1).
\item Communication and privacy settings must prevent interactions the user does not permit (FR-39.5, NFR-SEC-6).
\end{enumerate}
\par\medskip

\hypertarget{sec:b8-special-requirements}{}
\addcontentsline{toc}{subsubsection}{B8. Special Requirements}
\subsubsection*{B8. Special Requirements}
\par\medskip

\begingroup
\small
\renewcommand{\baselinestretch}{1.0}\selectfont
\renewcommand{\arraystretch}{1.15}
\setlength{\tabcolsep}{6pt}
\begin{longtable}{|>{\raggedright\arraybackslash}p{0.10\textwidth}|>{\raggedright\arraybackslash}p{0.83\textwidth}|}
\hline
\textbf{ID} & \textbf{Requirement} \\ \hline
\endhead
NFR-PRIV-1 & Users must be able to configure profile visibility using at least public and private levels. \\ \hline
NFR-PRIV-2 & Users must be able to allow or reject incoming friend requests. \\ \hline
NFR-PRIV-3 & Users must be able to control whether university, major, and academic year are visible. \\ \hline
NFR-PRIV-4 & Users must be able to show or hide their online status. \\ \hline
NFR-PRIV-5 & Users must be able to control access to documents and learning resources they share. \\ \hline
NFR-PRIV-6 & Users must be able to control whether other users can send messages or room invitations. \\ \hline
NFR-PRIV-7 & Privacy settings must apply consistently to posts, comments, member lists, and search results. \\ \hline
NFR-SEC-1 & Authentication is required before accessing protected user resources. \\ \hline
NFR-SEC-3 & Sensitive communication must be encrypted using TLS 1.2 or higher. \\ \hline
NFR-SEC-4 & Private forum and private-group content must be restricted to authorized users. \\ \hline
NFR-SEC-6 & Server-side authorization checks must prevent unauthorized access. \\ \hline
NFR-USE-4 & Primary interface actions must use clear labels and recognizable controls. \\ \hline
\end{longtable}
\endgroup
\par\medskip

\hypertarget{sec:b9-extension-points}{}
\addcontentsline{toc}{subsubsection}{B9. Extension Points}
\subsubsection*{B9. Extension Points}
\par\medskip

\begingroup
\small
\renewcommand{\baselinestretch}{1.0}\selectfont
\renewcommand{\arraystretch}{1.15}
\setlength{\tabcolsep}{6pt}
\begin{longtable}{|>{\raggedright\arraybackslash}p{0.44\textwidth}|>{\raggedright\arraybackslash}p{0.49\textwidth}|}
\hline
\textbf{Extension Point} & \textbf{Location in Flow of Events} \\ \hline
\endhead
Field-level visibility & Basic flow steps 2--3 \\ \hline
Search and status opt-outs & Basic flow step 3; alternatives A2 and A3 \\ \hline
Communication limits and member lists & Basic flow steps 4--5; alternative A1 \\ \hline
Global privacy enforcement & Basic flow step 8 \\ \hline
\end{longtable}
\endgroup
\par\medskip

\hypertarget{sec:b10-interface-and-data-notes}{}
\addcontentsline{toc}{subsubsection}{B10. Interface and Data Notes}
\subsubsection*{B10. Interface and Data Notes}
\par\medskip

\begin{itemize}
\item Screen/flow path: \texttt{Home → Settings → Privacy}; \texttt{Home → Settings → Communication} (recommended grouping in \texttt{settings.md}); profile settings flow \texttt{Edit Description → Privacy (Hidden | Visible)}, \texttt{Edit Hobbies → Privacy}, \texttt{Privacy → Submit}.
\item Exact values from \texttt{settings.md}: \textbf{Public / Private}, \textbf{Visible / Hidden}, \textbf{Everyone / Confirmed friends only}, \textbf{Allow / Reject incoming friend requests}, \textbf{Allow / Reject room invitations}, and \textbf{Applied to posts, comments, member lists, and search results}.
\item Exact buttons and actions: \textbf{Save Privacy}, \textbf{Save Communication}, \textbf{Mute}, \textbf{Unmute}, \textbf{Block}, \textbf{Unblock}, and the \texttt{ReportBehaviorLink} label \textbf{Report inappropriate behavior}.
\item Recommended components: \texttt{PrivacySettingsForm}, \texttt{ProfileVisibilitySelect}, \texttt{AcademicInfoVisibilitySelect}, \texttt{ContactInfoVisibilitySelect}, \texttt{StudyHistoryVisibilitySelect}, \texttt{HabitsVisibilitySelect}, \texttt{OnlineStatusToggle}, \texttt{LastActiveVisibilityToggle}, \texttt{PublicDiscoveryOptOutToggle}, \texttt{StudyPartnerSearchOptOutToggle}, \texttt{DirectMessageRestrictionSelect}, \texttt{FriendRequestPermissionSelect}, \texttt{RoomInvitationPermissionSelect}, \texttt{MutedUsersList}, \texttt{BlockedUsersList}.
\item Data fields: field-level \textbf{Visible / Hidden} values for bio, habits, study history, and contact info; \texttt{onlineStatusVisible}, \texttt{lastActiveVisible}; \texttt{directMessageRestriction}, \texttt{friendRequestsEnabled}, \texttt{roomInvitationsEnabled}, \texttt{mutedUserIds}, \texttt{blockedUserIds}.
\item Routes \texttt{/settings/privacy} and \texttt{/settings/communication} are recommendations in \texttt{settings.md}; confirmed routes are Not specified in proposals.
\item API: no dedicated privacy or communication endpoint, status code, or error code is documented in \texttt{API.md}.
\item Interaction rule: room invitation controls interact with FR-38.3 and AC-38.4 notification behavior, but this use case does not redefine notification settings.
\item Enforcement rule: hidden fields must not be exposed through badges, counters, or supporting text, and backend checks are required rather than UI-only hiding (NFR-PRIV-7, NFR-SEC-6).
\end{itemize}
\par\medskip

\hypertarget{sec:b11-source-alignment-notes}{}
\addcontentsline{toc}{subsubsection}{B11. Source Alignment Notes}
\subsubsection*{B11. Source Alignment Notes}
\par\medskip

\begin{itemize}
\item No dedicated privacy or communication API endpoints are documented in \texttt{API.md}; all route names in \texttt{settings.md} are recommendations.
\item \texttt{settings.md} treats Settings as an aggregate hub with no dedicated Settings use case; these controls are grouped from US-10 and US-39 without creating new use case IDs.
\item \texttt{friend\_profile.md} consumes these values to hide fields and restrict messaging, friend requests, and room invitations; it defines no API contract or error codes.
\item AC-39.1 and FR-39.3 mute/block management overlaps this use case only as settings control; request and message blocking flows remain in the friend and communication stories.
\end{itemize}
\par\medskip

\par\medskip
\hypertarget{sec:ui-prototype-specimens}{}
\addcontentsline{toc}{subsection}{UI Prototype Specimens}
\subsection*{UI Prototype Specimens}
\par\medskip
\begin{figure}[H]
\centering
\includegraphics[width=0.75\textwidth,height=0.42\textheight,keepaspectratio]{"../../prototype/profile_p6_private_profile.png"}
\par\vspace{2pt}
{\footnotesize \textbf{UI Prototype:} Profile: P6 Private profile}
\end{figure}
\par\smallskip
\begin{figure}[H]
\centering
\includegraphics[width=0.75\textwidth,height=0.42\textheight,keepaspectratio]{"../../prototype/settings_04_settings_privacy_and_notifications.png"}
\par\vspace{2pt}
{\footnotesize \textbf{UI Prototype:} Settings: 04 Settings - Privacy \& notifications}
\end{figure}
\par\smallskip
\begin{figure}[H]
\centering
\includegraphics[width=0.75\textwidth,height=0.42\textheight,keepaspectratio]{"../../prototype/settings_05_profile_validation_and_notification_selection.png"}
\par\vspace{2pt}
{\footnotesize \textbf{UI Prototype:} Settings: 05 - Profile validation \& notification selection}
\end{figure}
\par\smallskip

\end{document}
